% 5-3-pricecast.tex
%  Budget: ~2pp.  HANDOFF 12, 19.2 (figure 16). Code: app/pricecast.py,
%  app/forecast_service.py.
%  MANDATE:
%   - the tool demonstrates the pipeline; it is NOT a trading aid, and the
%     tool itself says so on screen -- that is the point worth writing
%   - every number here is already reported in 5-2 or visible in figure 16
%  Rules: decimals in \lr{}; integers Persian; every table cell in \lr{}.

\section{ابزار \lr{EnergyMarket-PriceCast}}
\label{sec:5-3}

خروجی کاربردیِ این پژوهش، یک ابزار تعاملی است که همان خط لوله‌ی فصل سوم را از
دریافت داده تا پیش‌بینیِ ۲۴ ساعتِ روز بعد، در یک رابط کاربری در دسترس می‌گذارد.
هدف از ساختن آن، نشان‌دادنِ اجراپذیریِ روش بود و نه ارائه‌ی یک سامانه‌ی معاملاتی؛
و همان‌طور که در ادامه می‌آید، خودِ ابزار این تمایز را به کاربرش اعلام می‌کند.

\subsection*{آنچه ابزار انجام می‌دهد}

کاربر یک منبع داده را از میان سه گزینه برمی‌گزیند: عکس فوریِ ذخیره‌شده که بدون
اتصال به شبکه کار می‌کند، فراخوانیِ زنده‌ی \lr{API} انرژی‌چارتس، یا بارگذاری یک
فایل \lr{CSV} با همان طرحواره. سپس روزِ هدف را از میان روزهای دارای تاریخچه‌ی
کافی انتخاب می‌کند. ابزار برای آن روز ۲۴ قیمت ساعتی تولید می‌کند، آن‌ها را در
کنار قیمت‌های محقق‌شده رسم می‌کند، و میانگین روزانه‌ی پیش‌بینی، میانگین روزانه‌ی
واقعی، \lr{MAE} همان روز و بدترین ساعت را گزارش می‌دهد. جدول کاملِ ۲۴ مقدار نیز
در یک بخش بازشونده در دسترس است.

\begin{figure}[htbp]
  \centering
  \includegraphics[width=0.95\textwidth]{16_pricecast_screenshot}
  \caption{ابزار \lr{EnergyMarket-PriceCast}: پیش‌بینی ساعتیِ روز بعد در برابر
  قیمت محقق‌شده، به‌همراه هشدارِ دقتی که خودِ ابزار نمایش می‌دهد.}
  \label{fig:pricecast}
\end{figure}

\subsection*{ساختار}

منطق پیش‌بینی از رابط کاربری جدا نگه داشته شده است. لایه‌ی سرویس، مدلِ
\lr{LightGBM} تثبیت‌شده را از پوشه‌ی مدل‌های منجمد بارگذاری می‌کند، پنجره‌ی
تاریخچه‌ی لازم را می‌سازد، و اگر تاریخچه برای ساختن ویژگی‌های تأخیری کافی نباشد
خطای صریح می‌دهد به‌جای آنکه ورودیِ ناقص را با مقدار جایگزین پر کند. همین جداییْ
معنایی بیش از نظم کد دارد: ابزار همان تابع ویژگی‌سازی و همان شیء مدلی را
فرامی‌خواند که نتایج فصل چهارم از آن آمده است، پس آنچه روی صفحه دیده می‌شود
بازپیاده‌سازیِ تقریبیِ روش نیست، خودِ آن است. انتساب داده‌ی انرژی‌چارتس نیز، همان
عبارتی که پروانه‌ی \lr{CC BY 4.0} اقتضا می‌کند، در تمام حالت‌ها در پایین صفحه
نمایش داده می‌شود.

\subsection*{هشداری که ابزار درباره‌ی خودش می‌دهد}

بالای هر پیش‌بینی، کادری ثابت این جمله را نشان می‌دهد که این پیش‌بینی‌ها بر
قیمت‌های امروز دقیق نیستند، و این یک یافته است نه یک ایراد. سپس همان اعداد
\cref{sec:5-2} را نقل می‌کند: مدل بر شرایط ۳۱ دسامبر ۲۰۱۷ تثبیت شده و بر داده‌ی
زنده‌ی ۲۰۲۶ \lr{rMAE} برابر \lr{1.83} می‌گیرد، یعنی بدتر از پیش‌بینیِ ساده با
\lr{0.81}. این اعداد در ابزار نوشته نشده‌اند، بلکه از فایل نتیجه‌ی تثبیت‌شده‌ی
همان آزمون خوانده می‌شوند؛ اگر آن نتیجه تغییر کند، متن هشدار نیز تغییر می‌کند و
امکان ندارد از آن عقب بماند.

نمونه‌ی \cref{fig:pricecast} همین را نشان می‌دهد: منحنی پیش‌بینی، شکل روزانه‌ی
قیمت را دنبال می‌کند و اوج عصرگاهی را در جای درست قرار می‌دهد، اما در تراز
به‌کلی پایین‌تر از قیمت محقق‌شده می‌ایستد \lr{—} همان جابه‌جایی سطحی که
\cref{sec:5-2} سازوکارش را شکافت، این بار به‌جای یک ستون در جدول، در یک نمودار.

قراردادنِ این هشدار در معرض دید، تصمیمی آگاهانه بود. ابزاری که پیش‌بینی‌های
نادرست را بی‌هیچ قیدی نمایش دهد، در بهترین حالت بی‌فایده و در بدترین حالت
گمراه‌کننده است. ابزاری که کیفیت خود را اعلام می‌کند، هم مصداقِ اجراپذیریِ روش
است و هم مصداقِ برخورد درست با یک نتیجه‌ی منفی.
